\par\Needspace{6\baselineskip}\medskip
\hypertarget{native-118-article088}{}\pdfbookmark[2]{88.}{article-088}
\begin{center}88.\end{center}
The exposition of the \textit{second} method will give us occasion to work out a great many new and elegant relations. Since these relations take different forms for the different kinds of conic section, each kind must be treated separately. We shall begin with the ELLIPSE.

\SourceLeaf{119}Let $v$ and $v'$ be the true anomalies at the two positions, with $v$ referring to the earlier position. Let $E$ and $E'$ be the corresponding eccentric anomalies, and $r$ and $r'$ the corresponding radius vectors. Further, let $p$ be the semiparameter, $e=\sin\varphi$ the eccentricity, $a$ the semimajor axis, and $t$ the time taken to move from the first position to the second. Finally, let us set $v'-v=2f$, $v'+v=2F$, $E'-E=2g$, $E'+E=2G$, and $a\cos\varphi=\dfrac{p}{\cos\varphi}=b$. Combining formulas V and VI of Art. 8 now readily gives the following equations.
\[
\begin{aligned}
\mathrm{1)}\qquad&b\sin g=\sin f\mathbin{.}\BareRoot{rr'}\\[4pt]
\mathrm{2)}\qquad&b\sin G=\sin F\mathbin{.}\BareRoot{rr'}
\end{aligned}
\]
\[
\begin{aligned}
p\cos g={}&\{\cos\tfrac12v\cos\tfrac12v'\mathbin{.}(1+e)\\[3pt]
&\quad+\sin\tfrac12v\sin\tfrac12v'\mathbin{.}(1-e)\}\BareRoot{rr'}
\end{aligned}
\]
or
\[
\mathrm{3)}\qquad p\cos g=(\cos f+e\cos F)\BareRoot{rr'}
\]
and similarly
\[
\mathrm{4)}\qquad p\cos G=(\cos F+e\cos f)\BareRoot{rr'}.
\]
Combining equations 3 and 4 then gives
\[
\begin{aligned}
\mathrm{5)}\qquad&\cos f\mathbin{.}\BareRoot{rr'}=(\cos g-e\cos G)a\\[4pt]
\mathrm{6)}\qquad&\cos F\mathbin{.}\BareRoot{rr'}=(\cos G-e\cos g)a.
\end{aligned}
\]
From formula III of Art. 8 we obtain
\[
\begin{aligned}
\mathrm{7)}\qquad r'-r&=2ae\sin g\sin G\\[4pt]
r'+r&=2a-2ae\cos g\cos G\\[4pt]
&=2a\sin g^2+2\cos f\cos g\mathbin{.}\BareRoot{rr'},
\end{aligned}
\]
hence
\[
\mathrm{8)}\qquad a=\frac{r+r'-2\cos f\cos g\mathbin{.}\BareRoot{rr'}}{2\sin g^2}\mathpunct{\cdot}
\]
Let us set
\[
\mathrm{9)}\qquad\frac{\BareRoot{\dfrac{r'}r}+\BareRoot{\dfrac{r}{r'}}}{2\cos f}=1+2l,
\]
then we have
\[
\mathrm{10)}\qquad a=\frac{2(l+\sin\tfrac12g^2)\cos f\mathbin{.}\BareRoot{rr'}}{\sin g^2}
\]
and also $\BareRoot{a}=\pm\dfrac{\BareRoot{(2(l+\sin\tfrac12g^2)\cos f\mathbin{.}\BareRoot{rr'})}}{\sin g}\mathpunct{\raisebox{.45ex}{,}}$ where the upper or lower sign is to be taken according as $\sin g$ is positive or negative. Formula XII of Art. 8 \SourceLeaf{120}gives us the equation
\[
\begin{aligned}
\frac{kt}{a^{3/2}}&=E'-e\sin E'-E+e\sin E\\[4pt]
&=2g-2e\sin g\cos G\\[4pt]
&=2g-\sin2g+2\cos f\sin g\frac{\BareRoot{rr'}}a\mathpunct{\cdot}
\end{aligned}
\]
If we now substitute for $a$ in this equation its value from equation 10 and, for brevity, set
\[
\mathrm{11)}\qquad\frac{kt}{2^{3/2}\cos f^{3/2}(rr')^{3/4}}=m,
\]
then, after simplifying, we obtain
\[
\mathrm{12)}\qquad\pm m=(l+\sin\tfrac12g^2)^{1/2}+(l+\sin\tfrac12g^2)^{3/2}\left(\frac{2g-\sin2g}{\sin g^3}\right),
\]
where the upper or lower sign is to be placed before $m$ according as $\sin g$ is positive or negative.

Whenever the heliocentric motion is between $180^0$ and $360^0$, or more generally whenever $\cos f$ is negative, the quantity $m$ determined by formula 11 would be imaginary and $l$ would be negative. To avoid this, we adopt the following equations in place of equations 9 and 11 in this case.
\[
\begin{aligned}
\mathrm{9^*)}\qquad&\frac{\BareRoot{\dfrac{r'}r}+\BareRoot{\dfrac{r}{r'}}}{2\cos f}=1-2L\\[8pt]
\mathrm{11^*)}\qquad&\frac{kt}{2^{3/2}(-\cos f)^{3/2}(rr')^{3/4}}=M,
\end{aligned}
\]
From these we obtain the following in place of equations 10 and 12.
\[
\begin{aligned}
\mathrm{10^*)}\qquad a&=\frac{-2(L-\sin\tfrac12g^2)\cos f\mathbin{.}\BareRoot{rr'}}{\sin g^2}\\[8pt]
\mathrm{12^*)}\qquad\pm M&=-(L-\sin\tfrac12g^2)^{1/2}+(L-\sin\tfrac12g^2)^{3/2}\left(\frac{2g-\sin2g}{\sin g^3}\right),
\end{aligned}
\]
Here the choice of sign is made in the same way as before.
